\subsection{Two-site exposure ordering}\label{app:two-site-example}

Whether a capacity response separates firms' business mixes depends on
their initial exposures.
To make that dependence explicit, consider two firms at two distinct activity
sites with identity demand loading and positive diagonal quadratic cost
curvature.
Initial interior capacities are $q_i=(q_{i1},q_{i2})$ and
$q_j=(q_{j1},q_{j2})$, all positive.
At activity 1, the node-1 opportunity value and outside appeal satisfy $r,b>0$,
and each firm's payoff component has the form
\begin{equation}\label{eq:sign-payoff}
  \pi_i^{(1)}(q_{i1},q_{j1})
  =\frac{rq_{i1}}{b+q_{i1}+q_{j1}}-c_i q_{i1}-\frac{\omega_i}{2}q_{i1}^2,
  \qquad \omega_i,\omega_j>0.
\end{equation}
Firm $j$ receives a small cost reduction $c_j\mapsto c_j-s$, preserving
nonnegative costs.
Assume a differentiable interior equilibrium path satisfying both first-order
conditions locally, with activity 2 unchanged.
To compare business mixes we use the quadratic Wasserstein distance $W_2$,
the least root-mean-square distance over which mass must move to turn one
composition into the other.
The two sites are a distance $d_*>0$ apart.
Define $P_k(s)$ from the normalized two-site capacities and
$Y^F_{ij}(s)=W_2^2(P_i(s),P_j(0))-W_2^2(P_i(0),P_j(0))$, the change in the
squared distance from firm $i$'s business mix to firm $j$'s initial mix.

\begin{proposition}[Conditional capacity and composition
  response]\label{prop:two-site-adjustment}
  In the two-site economy and along the path just specified, suppose
  \begin{equation}\label{eq:sign-conditions}
    q_{i1}<b+q_{j1},\qquad q_{j1}q_{i2}>q_{i1}q_{j2}.
  \end{equation}
  Then the cost recipient expands in activity 1, $\dot q_{j1}>0$, its competitor
  contracts, $\dot q_{i1}<0$, and $\dot Y^F_{ij}>0$.
  Reversing the strict initial exposure ordering reverses the sign of
  $\dot Y^F_{ij}$ while preserving the capacity-response signs.
\end{proposition}

\subsubsection{Proof of the two-site capacity and composition
response}\label{app:two-site-adjustment}

\begin{proof}[Proof of Proposition~\ref{prop:two-site-adjustment}]
  Differentiate the two interior first-order conditions in
  \eqref{eq:sign-payoff}.
  Write $D=b+q_{i1}+q_{j1}$, $B_i=b+q_{j1}$, $B_j=b+q_{i1}$, and define
  \begin{equation}\label{eq:sign-jacobian}
    \begin{aligned}
      J_{ii}&=\omega_i+2rB_i/D^3,& J_{jj}&=\omega_j+2rB_j/D^3,\\
      J_{ij}&=r(B_i-q_{i1})/D^3,& J_{ji}&=r(B_j-q_{j1})/D^3,\\
      \Delta&=J_{ii}J_{jj}-J_{ij}J_{ji}.
    \end{aligned}
  \end{equation}
  These are the own and rival entries of the field's
  Jacobian~\eqref{eq:jacobian} restricted to activity 1.
  The demand-only part of the determinant has numerator
  $r^2[4{B_i}B_j-(B_i-q_{i1})(B_j-q_{j1})]$ over $D^6$.
  The bracket equals $3b^2+4b(q_{i1}+q_{j1})+(q_{i1}+q_{j1})^2>0$.
  The additional terms from $\omega_i,\omega_j>0$ are positive, so $\Delta>0$.
  Solving the differentiated conditions gives
  \begin{equation}\label{eq:sign-capacity}
    \begin{pmatrix}J_{ii}&J_{ij}\\J_{ji}&J_{jj}
    \end{pmatrix}
    \begin{pmatrix}\dot q_{i1}\\\dot q_{j1}
    \end{pmatrix}
    =
    \begin{pmatrix}0\\1
    \end{pmatrix},
    \qquad \dot q_{i1}=-J_{ij}/\Delta,\quad \dot q_{j1}=J_{ii}/\Delta.
  \end{equation}
  The substitution inequality $q_{i1}<b+q_{j1}$ makes $J_{ij}>0$, yielding
  competitor contraction and recipient expansion.

  For shares $p,p'\in[0,1]$ on two sites separated by $d_*>0$, the exact
  squared transport distance is
  \begin{equation}\label{eq:two-site-transport}
    W_2^2\bigl((p,1-p),(p',1-p')\bigr)=d_*^2|p-p'|.
  \end{equation}
  Any feasible nonnegative coupling has off-diagonal masses whose
  difference is $p-p'$.
  Their sum is at least $|p-p'|$, and matching the common mass at each site
  attains the bound.
  Put $p_k=q_{k1}/(q_{k1}+q_{k2})$.
  Because activity 2 is unchanged,
  $\dot p_i=q_{i2}\dot q_{i1}/(q_{i1}+q_{i2})^2<0$.
  The exposure ordering implies $p_j(0)>p_i(0)$ and persists locally against
  the frozen $p_j(0)$.
  Differentiating the squared distance therefore gives
  \begin{equation}\label{eq:sign-distance}
    \dot Y^F_{ij}=-d_*^2\dot p_i
    =\frac{d_*^2{q_{i2}}J_{ij}}{(q_{i1}+q_{i2})^2\Delta}>0.
  \end{equation}
  If $p_j(0)<p_i(0)$, the derivative instead equals $d_*^2\dot p_i<0$.
  This establishes the statement and its reversed-order conclusion.
\end{proof}

Initial exposure ordering determines whether contraction moves the
competitor's business mix toward or away from the recipient's initial mix;
the capacity-response signs alone do not determine that direction.
